\chapter{Conclusion}\label{ch:conclusion}

This thesis built a digital-twin-based BMS with cloud monitoring. A twin of a 19S12P pack, fitted to measured M50T data, feeds a BMS controller on an STM32 over CAN, and a gateway forwards the telemetry to a web dashboard. The BMS estimates the SOC with 0.28\,\% RMSE, detects unexplained heating within seconds and limits the current in four graded states. It passed 1000 randomised runs on the microcontroller, and the cloud link delivered all telemetry without loss.

Compared with cloud monitoring systems such as \citet{nejakar2025}, the main step forward is control: the BMS does not only report a problem, it protects the battery, and the digital twin makes it possible to test this safely.

\paragraph{Future work.} Testing with real cells and a cell-monitoring board; capacity estimation; a learned SOC component combined with the EKF, as in the base paper; and a mobile interface.
